\documentclass[10pt,a4paper]{amsart}
\usepackage[margin=19mm]{geometry}
\usepackage{amsmath,amssymb,mathtools}
\usepackage[scr=rsfs,cal=euler]{mathalfa}
\usepackage{xfrac}
\usepackage[unicode,hidelinks,bookmarks=false]{hyperref}
\newcommand{\ie}{i.e.\ }
\newcommand{\cf}{cf.\ }
\newcommand{\resp}{respectively\ }
\newcommand{\Subsection}{\subsection}
\providecommand{\nobreakdash}{\nobreak\hbox{-}}
\setlength{\emergencystretch}{2em}
\multlinegap0pt
\newcommand{\ind}{{\mathrm{ind}}}
\usepackage{mathtools}
\usepackage{amssymb}
\usepackage{xcolor}
\usepackage{csquotes}
\usepackage{tikz}
\usepackage{float}
\newtheorem{lemma}{Lemma}
\newtheorem{remark}{Remark}
\newcommand{\De}{\Delta}
\newcommand{\cor}{\mathrm{cor}}
\newcommand{\cut}{\mathrm{cut}}
\renewcommand{\d}{\partial}
\newcommand{\de}{\delta}
\newcommand{\e}{\varepsilon}
\newcommand{\g}{\gamma}
\renewcommand{\k}{\kappa}
\newcommand{\s}{\sigma}
\newcommand{\w}{\omega}
\title{Sharp regularity of sub-Riemannian length-minimizing curves}
\author{Alessandro Socionovo}
\date{Source: arXiv:2502.00403v2 (CC BY 4.0)}
\begin{document}
\maketitle

\noindent\emph{Source and licence.} Sharp regularity of sub-Riemannian length-minimizing curves. Version arXiv:2502.00403v2 (CC BY 4.0), distributed by arXiv under the Creative Commons Attribution 4.0 International licence, http://creativecommons.org/licenses/by/4.0/, as recorded in the arXiv licence metadata for this exact version and shown on the abstract page https://arxiv.org/abs/2502.00403v2. Full licence notice: \texttt{licenses/arxiv-2502.00403-CC-BY-4.0.txt}.\par\smallskip

\noindent\textit{Editorial scope.} Counted unit: the article's lemma lem:correction (source0.tex lines 1176--1183). The article's complete original proof of it is delivered with it (source0.tex lines 1185--1241, one \texttt{proof} environment of 57 lines). No mathematics is written, repaired or re-derived: the statement and the proof are the article's own lines, in the article's own notation, and no retained line is substituted or re-flowed.  Retained uncounted as the source-local context the proof needs: lines 981--989; lines 1166--1169.  Nothing was omitted from any retained span, so the package carries no omission record.  No retained line contains a citation, so the wrapper carries no bibliography and no bibliography entry.  No retained line contains a figure environment, an image-inclusion command or drawing code, so no image is delivered and the package declares no asset dependency.  Editorial formatting only, in addition: an AMS \texttt{amsart} preamble that restates the article's own declarations and macros used by the retained lines, namely the environments \texttt{lemma}, \texttt{remark} on separate counters, together with the standard packages those lines need.  The retained environments print in this document as Lemma 1, Remark 1 in order of appearance, whereas the article prints the same items with the numbering of its own full document.  The complete native source of the article as distributed in the arXiv e-print for this exact version is bundled unchanged as \texttt{sources/172571/source0.tex}.
\begin{lemma}
	\label{lem:Stokes}
	For any $\omega \in \mathcal C_\e([0,\tau],\mathbb R^2)$ we have 
	\begin{equation}
		\label{eq:Stokes}
		\sum_{U} \mathscr L_Q(U)=0,
	\end{equation} 
	where the sum is over all connected components $U\subset\Omega$.
\end{lemma}

\begin{remark}
	\label{rem:win_num}
	The region enclosed by the curve $\w_{\rho,\de}*(-\g_\e)$ has two connected components, which are the rectangle $E_\de$ and the region enclosed by $\k_\rho * (-\g_\rho)$, which we call $U_{\rho}$. The winding numbers of $U_\rho$ and $E_\de$ have absolute value equal to 1, and we also have $Q|_{U_\rho}, Q|_{E_\de}\geq0$. Then the parameterization $\s_\de$ of $\d E_\de$ has been chosen in such a way that $\ind(E_\de)=-\ind(U_\rho)$, so that~\eqref{eq:Stokes} can be satisfied and $\w_{\rho,\de}$ may be a competitor.
\end{remark}

\begin{lemma}
	\label{lem:correction}
	For every $\e>0$ and $\rho,\de\in(0,\e)$ small enough there exist a constant $C=C(a,b,\e)>0$ and a function $f_{a,b}(\de)$, with $f_{a,b}(\de)=O(\de)$ as $\de\to0$, such that the curve $\w_{\rho,\de}$ is a competitor as soon as
	\begin{equation}
		\label{eq:correction}
		\rho^{2b+1}=C\de^3(1+f_{a,b}(\de)).
	\end{equation}
\end{lemma}

\begin{proof}
	We have to compute and then compare the effects made by the cut and by the rectangle on the third coordinate, as $\rho,\de\to0$.
	
	For the cut we have
	\begin{equation}
		\label{eq:errk}
		\begin{split}
			\De^{\cut}_3(\rho)&:=\int_{0}^\rho P(\k(t))^2\dot\k_2(t)dt=\int_{0}^{\rho} (t^a\rho^{b-a}-t^b)^2dt\\
			&\stackrel{(t=\rho s)}{=}\rho^{2b+1}\int_{0}^{1} (s^a-s^b)^2ds.
		\end{split}
	\end{equation}

	For the rectangle we use the Stokes theorem, see Lemma~\ref{lem:Stokes}. We compute
	\begin{equation}
		\label{eq:errQ'}
		\begin{split}	
			\De_3^{\cor}(\de)&:=\iint_{E_\de} Q(x_1,x_2)dx_1dx_2\\
			&=\int_{\e(1-\de)}^\e\int_{\e^q}^{\e^q(1+\de)}2ax_1^{a-1}(x_1^a-x_2^b)dx_1dx_2 \qquad (x_1=\e^qx,\; x_2=\e y)\\
			&=2a\e^{2b+1}\left(\int_{1-\de}^1\int_{1}^{1+\de} x^{a-1}(x^a-y^b)dxdy\right)
			%&=a(b+a)\e^{2b+1}\de^3(1+f_{a,b}(\de)),
		\end{split}
	\end{equation}
    We compute the integral appearing in the last line. We have
    \begin{equation}
        \label{eq:calcoloint1}
        \begin{split}
            &\int_{1-\de}^1\int_{1}^{1+\de} x^{a-1}(x^a-y^b)dxdy 
            \\
            &= \de\frac{(1+\de)^2a-1}{2a} - \frac{(1+\de)^a-1}{a}\frac{1-(1-\de)^b}{b+1}
            \\
            &=\frac{1}{2a} \sum_{k=1}^{2a}\binom{2a}{k}\de^{k+1} - 
            \frac{1}{a(b+1)} \sum_{i=1}^{a}\sum_{j=1}^{b+1}(-1)^{j+1}\binom{a}{i}\binom{b+1}{j}\de^{i+j}
            \\
            &:= S_1 - S_2
        \end{split}
    \end{equation}
    The two sums $S_1$ and $S_2$ above reads in the following way
    \begin{align}
        \label{eq:S1}&S_1 = \de^2 + \frac{2a-1}{2}\de^3 + \frac{1}{2a} \sum_{k=3}^{2a}\binom{2a}{k}\de^{k+1}
        \\
        \label{eq:S2}&S_2 = \de^2 - \frac b2\de^3 + \frac{a-1}{2} \de^3 + \frac{1}{a(b+1)} \sum_{i=2}^{a}\sum_{j=2}^{b+1}(-1)^{j+1}\binom{a}{i}\binom{b+1}{j}\de^{i+j}
    \end{align}
    From \eqref{eq:errQ'}, \eqref{eq:calcoloint1}, \eqref{eq:S1}, and \eqref{eq:S2}, we deduce
    \begin{equation}
        \label{eq:errQ}
        \De_3^{\cor}(\de)=a(a+b)\e^{2b+1}\de^3(1+f_{a,b}(\de)),
    \end{equation}
    with
    \begin{equation}
        \label{eq:fabdelta}
        f_{a,b}(\de)= \frac{1}{2a} \sum_{k=3}^{2a}\binom{2a}{k}\de^{k-2} - \frac{1}{a(b+1)} \sum_{i=2}^{a}\sum_{j=2}^{b+1}(-1)^{j+1}\binom{a}{i}\binom{b+1}{j}\de^{i+j-3}.
    \end{equation}
    
    
	
	By Remark~\ref{rem:win_num}, the effects on the third coordinate produced by the cut and by the rectangle have opposite sign. Therefore, the proof of~\eqref{eq:correction} follows by~\eqref{eq:errk} and~\eqref{eq:errQ} by taking $C=C(a,b,\e)=a(b+a)\e^{2b+1}\big(\int_{0}^{1} (s^a-s^b)^2ds\big)^{-1}>0$ and $\rho,\de>0$ small enough. Moreover, by \eqref{eq:fabdelta} it follows $f_{a,b}(\de)=O(\de)$, as $\de\to0$.
\end{proof}

\end{document}
